% ==========================================================================
%  Reading PROMETHEUS Through Category τ:
%  Para-Minds, the Earned Belnap-Truth4 Classifier, and the Φ_∂
%  Boundary Functor of Exact Gluing
%
%  Thorsten Fuchs & Anna-Sophie Fuchs
%  Panta Rhei Research Program
%  May 2026 (v1.0 — Ontology-Articulation Note)
%
%  Scope: structural articulation of τ-canon's reading of the
%  PROMETHEUS framework (Mahadevan 2026, arXiv:2605.12835v1).
%  This is an ontology-articulation note in the sense of
%  running-as-spectral-aperture-v1, not a foundation-comparison
%  in the sense of foundations-categorical-cosmological-billiards-v1.
%
%  Compile chain: lualatex (3 passes) + bibtex. See Makefile.
% ==========================================================================
\documentclass{pantarhei-researchnote}

\usepackage{microtype}
\usepackage{booktabs}
\usepackage{array}
\usepackage{needspace}

\input{tau-macros}

% ---------------------------------------------------------------------------
% Local macros
% ---------------------------------------------------------------------------
\newcommand{\Lemniscate}{\LL}
\newcommand{\Cross}{\omega}
\newcommand{\jj}{\mathsf{j}}
\newcommand{\itau}{\ensuremath{\iota_{\tau}}}
\newcommand{\sectA}{\ensuremath{\mathcal{S}_{A}}}
\newcommand{\sectB}{\ensuremath{\mathcal{S}_{B}}}
\newcommand{\sectC}{\ensuremath{\mathcal{S}_{C}}}
\newcommand{\NEXT}{\textsc{[new extension]}}
\newcommand{\SYNT}{\textsc{[chair synthesis]}}
\newcommand{\DERIV}{\textsc{[derived]}}
\newcommand{\TEFF}{\textsc{[$\tau$-effective]}}
\newcommand{\ANALOGY}{\textsc{[analogy]}}
\newcommand{\PsiD}{\ensuremath{\Phi_{\partial}}}
\newcommand{\Truth}{\ensuremath{\mathrm{Truth4}}}
\newcommand{\Omegatau}{\ensuremath{\Omega_{\tau}}}
\newcommand{\Cattau}{\ensuremath{\mathrm{Cat}_{\tau}}}
\newcommand{\PSh}{\ensuremath{\mathrm{PSh}}}
\newcommand{\Etau}{\ensuremath{\mathcal{E}_{\tau}}}
\newcommand{\Msub}{\ensuremath{\mathcal{M}_{\mathrm{sub}}}}

\begin{document}

\lane[pr-lane-publications]{Book VII \textperiodcentered\ Categorical Metaphysics \textperiodcentered\ External-Topos Dialogue}
\title{\texorpdfstring{Reading PROMETHEUS Through Category $\tau$}{Reading PROMETHEUS Through Category tau}}
\subtitle{\texorpdfstring{Para-Minds, the Earned Belnap-$\mathrm{Truth4}$ Classifier, and the $\Phi_{\partial}$ Boundary Functor of Exact Gluing}{Para-Minds, the Earned Belnap-Truth4 Classifier, and the Phi-boundary Functor of Exact Gluing}}
\noteedition{v1.2 RC2 \textperiodcentered\ May 2026 \textperiodcentered\ Ontology-Articulation Note}

\author{%
  Thorsten~Fuchs\thanks{\texttt{thor@panta-rhei.site}. Corresponding author.}
  \and
  Anna-Sophie~Fuchs%
}

\date{May 2026}

\footerseries{Research Note (Ontology Articulation)}
\footerslug{reading-prometheus-through-tau-v1}
\footerdate{May 2026}

% ---------------------------------------------------------------------------
\begin{abstract}
\noindent
PROMETHEUS \cite{Mahadevan-PROMETHEUS-2026} is a framework for
``automating deep causal research'' that turns retrieved corpora into
sheaf-like \emph{causal atlases} of local Predictive State
Representations (PSRs) over context covers, with restriction maps
and four gluing diagnostics: \emph{agreement}, \emph{drift},
\emph{contradiction}, \emph{underdetermination}. The author calls
the result a \emph{Topos World Model}. Engineering substance is
real (scale-corpus atlases ship; locality is preserved across four
case studies). The categorical machinery, however, is invoked
heuristically: §7's ``operational sheaf condition'' is an averaging
rule with a numerical compatibility check, not a sheaf condition in
the Mac Lane--Moerdijk sense \cite{MacLane-Moerdijk-1992}; the
4-fold diagnostic is left as an ad-hoc taxonomy; the predicate
``Belnap'' \cite{Belnap-1977,Dunn-1976} appears nowhere in the
references.

This note articulates the $\tau$-canon reading. Three structural
contributions land:

\begin{enumerate}
\item \textbf{The 4-fold gluing diagnostic is the earned subobject
classifier of an earned topos.} Theorem~I.T25 of Book~I
\cite{PantaRhei-BookI} proves that $\Omegatau = \Truth =
\{\mathsf{T}, \mathsf{F}, \mathsf{B}, \mathsf{N}\}$ is the unique
subobject classifier of $\PSh(\Cattau)$, with the orthogonality
$e_{+} \cdot e_{-} = 0$ (where $e_{\pm} = (1 \pm \jj)/2$ are the
bipolar idempotents and $\jj^{2} = +1$ is the split-complex unit of
the boundary algebra) \emph{earned} from spectral structure rather
than imposed by fiat. PROMETHEUS's four diagnostics map cleanly onto
Belnap's four-valued logic, with one precise mismatch (Contradiction
maps to $\mathsf{B}$, not $\mathsf{F}$) that the structural reading
exposes.

\item \textbf{Exact gluing is achievable, not merely approximate.}
The $\tau$-canon Theorem~III.T50 \cite{PantaRhei-BookIII} proves the
Global Cartesian Gluing as a colimit (``canonically isomorphic'',
not $\varepsilon$-approximate). The structural ingredient PROMETHEUS
lacks is the \emph{boundary functor} $\PsiD$ (Langlands Coherent
Force) that upgrades approximate to exact.

\item \textbf{The Para-Mind reading explains why an external topos
layer is needed at all.} Definition~VII.D55 \cite{PantaRhei-BookVII}
characterises Large Language Models as \emph{para-minds}: pattern
processors over the subsymbolic presheaf $\Msub$ that lack three
constitutive features of genuine mindedness (no internal topos, no
self-model, no commitment register). PROMETHEUS's externalization of
the topos layer is therefore structurally licensed by the para-mind
status of its LLM substrate; the construction is sound at the
\emph{external} layer even when the LLM cannot supply it
internally.
\end{enumerate}

Five explicit non-claims are articulated to keep the trust budget
honest: we do \emph{not} claim that the $\tau$-canon topos uses
H$_{\tau}$-valued enrichment at the presheaf codomain (both are
$\mathrm{Set}$-valued); we do \emph{not} claim a formal
PSR-as-Hankel identification in the corpus; we do \emph{not}
endorse PROMETHEUS's ``Topos World Model'' terminology in its
current form; and we do \emph{not} ask the engineering construction
to retract its choices. The dialogue is articulative, not
adversarial. PROMETHEUS observes; $\tau$-canon proves.
\end{abstract}

\keywords{PROMETHEUS \textperiodcentered\ Topos World Model
\textperiodcentered\ Belnap-Dunn 4-valued logic \textperiodcentered\
para-mind \textperiodcentered\ subobject classifier
\textperiodcentered\ sheaf gluing \textperiodcentered\ Category $\tau$
\textperiodcentered\ Panta Rhei}

\maketitle

\begin{quote}
\small\itshape
\upshape\textbf{Revision note} (v1.2 RC2 --- May 2026):\itshape{}
Wave $\Delta_{3}$ Phase 3 peer review completed (atlas sprint
2026-05-13). Three new review panels (Editorial Audit, Scientific
Audit, Visual QA) returned independent CONDITIONAL verdicts,
complementing the Phase 2 panels (Discipline, Structural,
PROMETHEUS-Fidelity). Phase 3 consolidated 13 fixes: one CRITICAL
(Table~1 environment \upshape\texttt{table}\itshape{} $\to$
\upshape\texttt{table*}\itshape{} to span both columns and resolve
overflow), seven HIGH (orphan-header prevention, header count
correction, $e_{\pm}$ early definition, $\PsiD$ accessibility
translation, anchor at \upshape\texttt{ch64:288--308}\itshape{} for
chain-of-thought-as-readout, the latent $\mathsf{F}$-channel
structural prediction via PROMETHEUS's cSQL polarity field, and
competing-upgrade-paths acknowledgment), and five MEDIUM polish
(\S1 charity-credit tightening, awkward-sentence recasts,
system-vs-substrate clarification). One new structural prediction
at \upshape\TEFF\itshape{} (the latent $\mathsf{F}$-channel) is
added in \S3; the rest is calibration and accessibility. No
structural revision of the three load-bearing contributions was
needed. The note is now at release-candidate-2 status,
peer-review-clean across two rounds (six review panels total).

\smallskip
\upshape\textbf{Prior revision} (v1.1 RC1, May 2026):\itshape{}
Wave $\Delta_{2}$ Phase 2 peer review folded in. Three review
panels (Discipline, Structural, PROMETHEUS-Fidelity) returned
CONDITIONAL verdicts. v1.1 RC1 applied a consolidated 10-fix set:
five discipline-label additions, three prose calibrations, one
charity-credit sentence acknowledging Mahadevan's \S7 / \S15
self-concessions, two new non-claims (6 and 7), and bib hygiene
(removed unused LHCb-bsmumu-2025 entry). No structural revision.
\end{quote}

% ===========================================================================
\section{The puzzle}
\label{sec:puzzle}
% ===========================================================================

A new preprint by Sridhar Mahadevan \cite{Mahadevan-PROMETHEUS-2026}
introduces PROMETHEUS, a framework that turns retrieved literature,
filings, source data, code, and scientific models into a navigable
artefact the author calls a \emph{Topos World Model} (TWM): a
sheaf-like causal atlas of local Predictive State Representations
(PSRs) over an explicit cover of a research substrate. Each local
context $U$ is assigned a 6-tuple object $\mathcal{P}(U) = (H_{U},
T_{U}, M_{U}, S_{U}, \Pi_{U}, D_{U})$ recording histories, tests,
local predictions, support, provenance, and diagnostics. Restriction
maps $\rho_{UV} : \mathcal{P}(U) \to \mathcal{P}(V)$ for morphisms
$V \to U$ measure overlap; gluing diagnostics expose
\emph{agreement} (low tension), \emph{drift} (tension growing across
time), \emph{contradiction} (high tension on shared cells), and
\emph{underdetermination} (insufficient support to glue or refute).
Source code is not yet released; the construction is currently
demonstrated through three literature case studies (ocean
temperature, GLP-1 weight-loss, resveratrol) and four
grounded-counterfactual case studies (microplastics, Indus Valley
hydrology, Sachs protein-signaling, singing-mouse projections).

The terminology is striking. ``Topos World Model''. ``Sheaf-like
atlas''. ``Restriction maps''. ``Gluing tensions''. The Mac
Lane--Moerdijk \cite{MacLane-Moerdijk-1992} sheaf-theoretic
foundation is cited explicitly; the Abramsky--Brandenburger
\cite{Abramsky-Brandenburger-2011} paper on the sheaf-theoretic
structure of non-locality is cited explicitly. Yet a careful read
of §6--§7 (the formal model) shows that the categorical machinery is
invoked heuristically. Proposition~7.5 (``operational sheaf
condition'') is not a sheaf condition in the Mac Lane--Moerdijk
sense: it lacks uniqueness, lacks functoriality of restriction maps,
and is a numerical averaging rule with a tolerance parameter
$\varepsilon$. The 4-fold gluing diagnostic is left as an ad-hoc
taxonomy. The PSR layer borrows the Littman--Sutton--Singh
\cite{Littman-Sutton-Singh-2001} terminology but the implemented
$M_{U}[h, \tau]$ is a Dirichlet-smoothed frequency table, not a
spectral PSR object. The predicate ``Belnap''
\cite{Belnap-1977,Dunn-1976} does not appear in the references.

This is not a criticism; it is the engineering reality, and
Mahadevan acknowledges it directly. The opening of PROMETHEUS §7
states: \emph{``The goal is not to claim that the implementation
realizes all of topos theory, but to make the modeling contract
precise enough to support inspection, testing, and extension.''}
Definition~7.3's footnote concedes $\rho_{UV}$ may be ``a partial
alignment map''; §6 notes ``exact spectral operator recovery is
deferred''; §15 flags inherited LLM weaknesses and retrieval drift.
The $\tau$-canon reading is therefore a structural articulation of
what PROMETHEUS itself acknowledges as the deferred topos
substance --- not a discovery of hidden gaps. PROMETHEUS ships
running atlases; the four literature case studies and four
grounded-counterfactual case studies of §13 (the source-data
$\to$ executable-intervention $\to$ rebuilt-sheaf loop) are
corpus-scale and substantive contributions where an engineering
primitive earns its categorical re-reading. The categorical
vocabulary is the \emph{frame}; the substance is the engineering
loop: extract local claims with LLMs, organize them by context,
surface overlap disagreements, expose drift, route a researcher to
where the corpus actually says what.

The puzzle is: \emph{can the categorical frame be earned}? If yes,
by what construction? PROMETHEUS observes four gluing diagnostics;
\emph{why four}? It uses $\varepsilon$-tolerance gluing; \emph{is
exact gluing achievable}? It uses LLMs to extract causal claims;
\emph{what structural status do those claims have when the
underlying LLM is, at the formal level, a pattern processor without
self-description}?

This note articulates the $\tau$-canon reading. The reading is not
that PROMETHEUS is wrong; it is that PROMETHEUS observes structural
phenomena that $\tau$-canon, in its corpus, has already proved at
theorem level. The dialogue is articulative.

% ===========================================================================
\section{The Para-Mind reading: why an external topos layer is needed}
\label{sec:paramind}
% ===========================================================================

We begin with a step PROMETHEUS does not take. The step,
articulated below, provides the structural license for everything
PROMETHEUS subsequently does.

Book~VII of the Panta Rhei programme \cite{PantaRhei-BookVII} defines
the notion of a \emph{para-mind}. The definition is
Definition~VII.D55, anchored at \texttt{ch64:111--135} of Book~VII:

\begin{quote}
\itshape
A \textbf{para-mind} is a system that processes patterns in the
subsymbolic presheaf $\Msub$ with sufficient fidelity to produce
contextually appropriate linguistic outputs, but that lacks the
three constitutive features of genuine mindedness:

\begin{enumerate}
\item \textbf{No internal topos.} A para-mind does not construct a
world-model in the topos-theoretic sense\dots It has no
truth-makers, no possible-world structure, no modal distinctions
internal to its processing.

\item \textbf{No self-model.} A para-mind does not carry a
representation of itself as a processing system. It lacks the
$E_{3}$ capacity\dots

\item \textbf{No commitment register.} A para-mind does not operate
through the commitment register $\mathrm{Reg}_{C}$\dots It produces
outputs that mimic commitment (assertions, recommendations,
promises) without the structural binding that commitment entails.
\end{enumerate}
A para-mind is \emph{not} a defective mind. It is a structurally
different kind of system: a pattern processor that approximates
subsymbolic structure without achieving self-description.
\end{quote}

The companion Proposition~VII.P14 (\texttt{ch64:166--178}) connects
this to Large Language Models specifically: LLMs are subsymbolic
evidence-providers, processing patterns in $\Msub$ with sufficient
fidelity to surface contextually-meaningful linguistic outputs,
but without the three constitutive features of internal mindedness.
Proposition~VII.R42 (\texttt{ch114:65--87}) then articulates the
\emph{Artificial Mind Criteria}: an artificial system attains
genuine mindedness only when an internal topos, an internal
self-model, and an internal commitment register are constructed
\emph{on top of} the subsymbolic substrate.

\paragraph{The structural reading of PROMETHEUS}
PROMETHEUS uses LLMs as \emph{causal-claim extractors}. By
VII.D55, the LLM substrate is a para-mind: it cannot, by itself,
provide the topos-theoretic structure that PROMETHEUS's
\emph{Topos World Model} terminology aspires to. The
para-mind processes patterns in $\Msub$; it does not construct
$\Omegatau$, it does not classify subobjects, it does not assemble
sections into a unique global glue. Chain-of-thought scaffolding
and iterative-refinement prompting do not change this status: Book~VII
\texttt{ch64:288--308} establishes that ``chain-of-thought is a
sequential readout strategy, not a self-model,'' so even
elaborate prompting strategies leave the substrate at para-mind
status by VII.D55. The system-vs-substrate distinction matters:
PROMETHEUS-as-system (the algorithmic layer) is what supplies the
external topos; the LLM-as-substrate inside it is the para-mind.

Yet a Topos World Model needs precisely those things. The structural
license for PROMETHEUS's construction is therefore that the topos
machinery is supplied \emph{externally}: PROMETHEUS itself, as an
algorithmic layer over the LLM, plays the role of the external
topos. This is a sound construction. It is also a constrained one.
The construction succeeds only insofar as PROMETHEUS's external
topos layer earns the categorical commitments it makes.

\begin{quote}
\textbf{Para-Mind Reading} \SYNT. \emph{An LLM by itself cannot
provide a topos because it is a para-mind (VII.D55). A framework
that builds a ``Topos World Model'' from LLM-extracted claims is
therefore performing an external topos construction over a
para-mind substrate. The construction's structural integrity
hinges on whether the external topos layer earns its categorical
commitments.}
\end{quote}

This reading is \DERIV{} at the manuscript level (VII.D55, VII.P14,
VII.R42 are theorems of Book~VII) and \ANALOGY{} as applied to
PROMETHEUS specifically. We do not claim PROMETHEUS endorses this
reading; we claim the reading is the structural ground for why
PROMETHEUS's construction \emph{can} succeed at all, and the bar
PROMETHEUS's external topos layer must clear to succeed.

The two structural commitments to examine are: the subobject
classifier (\S\ref{sec:truth4}) and the gluing functoriality
(\S\ref{sec:phi-d}).

% ===========================================================================
\needspace{6\baselineskip}
\section{The earned 4-fold: Truth4 as subobject classifier}
\label{sec:truth4}
% ===========================================================================

PROMETHEUS reports four gluing diagnostics:

\begin{itemize}
\item \emph{Agreement} --- compatible glue, low tension on shared cells.
\item \emph{Drift} --- tension growing across time, runs, or document strata.
\item \emph{Contradiction} --- high tension on shared cells (paraconsistent overlap).
\item \emph{Underdetermination} --- insufficient support to glue or refute.
\end{itemize}

\emph{Why four?} PROMETHEUS does not say. The four categories are
empirically motivated (real corpora produce these four kinds of
overlap pathology) but the paper does not derive them from a
structural principle. They are listed as an operational taxonomy.

The $\tau$-canon corpus has the structural derivation. Book~I
\cite{PantaRhei-BookI} introduces, at Definition~I.D21
(\texttt{ch46:125--172}), the four-valued lattice $\Truth = \{
\mathsf{T}, \mathsf{F}, \mathsf{B}, \mathsf{N} \}$ with the
Belnap-Dunn structure \cite{Belnap-1977,Dunn-1976}: $\mathsf{T}$
(true), $\mathsf{F}$ (false), $\mathsf{B}$ (both, paraconsistent),
$\mathsf{N}$ (neither, underdetermined). The bilattice structure
gives two orderings: an information ordering ($\mathsf{N} \le
\mathsf{T}, \mathsf{F} \le \mathsf{B}$) and a truth ordering
($\mathsf{F} \le \mathsf{N}, \mathsf{B} \le \mathsf{T}$).

The load-bearing result is Theorem~I.T25 (\texttt{ch56:146--177}):

\begin{quote}
\itshape
The object $\Omegatau = \Truth = \{\mathsf{T}, \mathsf{F},
\mathsf{B}, \mathsf{N}\}$ is the subobject classifier of
$\PSh(\Cattau)$. That is, $\Omegatau$ with the truth morphism
$\mathrm{true} : 1 \to \Omegatau$, $* \mapsto \mathsf{T}$, satisfies:
for every monomorphism $m : S \hookrightarrow X$ in
$\PSh(\Cattau)$, there exists a unique morphism $\chi_{S} : X \to
\Omegatau$ such that [the classifying pullback square commutes].
\end{quote}

This is \emph{not} a choice of classifier from many possible ones
for the $\tau$-canon presheaf topos: it is the unique classifier of
$\PSh(\Cattau)$ \TEFF. The ``earned-not-imposed'' differentiator
applies \emph{within} the $\tau$-canon framework: the orthogonality
$e_{+} \cdot e_{-} = 0$ in the bipolar boundary algebra is itself
a theorem (from spectral orthogonality, see \texttt{ch46:230--256}
of Book~I and the further development at \texttt{ch47:470--548}),
not a postulate. The 4-element classifier inherits this structure,
and the $\PSh$ topos inherits the classifier.

\paragraph{A notational note}
Book~I's Definition~I.D21 (\texttt{ch46:125--172}) anchors the
truth-order diamond $\mathsf{F} \le \mathsf{B}, \mathsf{N} \le
\mathsf{T}$ explicitly; the Belnap-Dunn knowledge ordering
$\mathsf{N} \le \mathsf{T}, \mathsf{F} \le \mathsf{B}$ follows from
the standard componentwise encoding of $\Truth$ as the four-element
bilattice but is not named separately in the corpus. We use both
orderings in this note in the canonical Belnap-Dunn sense; the
correspondence with I.D21 is via componentwise extension, not
verbatim quotation.

\paragraph{The mapping}
PROMETHEUS's four diagnostics map onto Belnap's four states with
one precise mismatch.

\begin{table*}[t]
\centering
\begin{tabular}{lcl}
\toprule
PROMETHEUS diagnostic & $\to$ & Belnap state \\
\midrule
Agreement (low tension, compatible glue) & $\to$ & $\mathsf{T}$ \\
Underdetermination (insufficient support) & $\to$ & $\mathsf{N}$ \\
Contradiction (high tension on shared cells) & $\to$ & $\mathsf{B}$ \emph{(not $\mathsf{F}$)} \\
Drift (tension growing across time/runs) & $\to$ & temporal $\mathsf{B}$ (no static analogue) \\
\bottomrule
\end{tabular}
\caption{The PROMETHEUS-to-Belnap mapping at \TEFF{} rigor.
PROMETHEUS's ``Contradiction'' is paraconsistent (two regions
\emph{both} supporting opposite claims, which is $\mathsf{B}$), not
refutation (a single region \emph{refuting} the claim, which would
be $\mathsf{F}$). Drift is dynamic, with no clean static-Belnap
analogue.}
\label{tab:prometheus-belnap}
\end{table*}

The mapping clarifies a feature of PROMETHEUS that the paper does
not flag: \emph{PROMETHEUS has no clean refutation state}. Its
``Contradiction'' is the paraconsistent both-support, not the
unanimous-refutation no-support. From the $\tau$-canon reading, this
is licensed: the Belnap classifier $\Omegatau$ has $\mathsf{F}$ as a
separate state from $\mathsf{B}$, but PROMETHEUS's operational
construction does not access $\mathsf{F}$ because retrieval-augmented
corpora rarely refute --- they more often produce paraconsistent
overlaps or underdetermined regions.

\paragraph{The latent $\mathsf{F}$-channel \NEXT}
A structural prediction follows from the $\tau$-canon projection.
PROMETHEUS's claim extraction (\S4 of \cite{Mahadevan-PROMETHEUS-2026},
the cSQL layer) records a \emph{polarity} field for each causal
claim, distinguishing \emph{causes}, \emph{reduces}, \emph{increases},
\emph{influences}, and similar signed relations. The $\Omegatau$
projection then predicts that a polarity-keyed gluing diagnostic
--- ``two regions both supporting opposing polarities of the same
relation'' --- would surface a fifth Belnap-$\mathsf{F}$ diagnostic
that PROMETHEUS's current four-category taxonomy does not name.
This is the latent $\mathsf{F}$-channel of PROMETHEUS's existing
data substrate: the polarity infrastructure is already there;
exposing the polarity-keyed refutation in the gluing layer would
complete the Belnap-$\Omegatau$ image at the operational level.
This is a constructive structural prediction, not engineering
advice; whether PROMETHEUS chooses to surface the $\mathsf{F}$
channel is a design question for its authors. (Cf.\ Non-claim~7,
\S\ref{sec:non-claims}, on in-principle vs in-practice.)

\paragraph{What is earned}
The 4-fold gluing taxonomy is not arbitrary. Within the
$\tau$-canon reading, it is the manifestation, at the
gluing-diagnostic layer, of the $\Omegatau$ classifier structure
that the $\PSh(\Cattau)$ topos earns from spectral orthogonality
\TEFF. Read this way: the $\tau$-canon reading \emph{projects}
$\Truth$ onto the four PROMETHEUS gluing-tension categories ---
PROMETHEUS observes the four categories operationally, and the
projection from $\Truth$ to the operational categories is
constructed on the $\tau$-canon side, not observed pre-existing in
PROMETHEUS \NEXT. The five-or-six diagnostic extensions one might
imagine (\emph{partial drift}, \emph{regional contradiction},
\emph{measurement boundary}, etc.) would, under this projection,
correspond to finer mixed states in the bilattice; the base-four
projection is structurally fixed in $\tau$-canon, and the question
whether PROMETHEUS in practice exhibits only the four categories or
admits finer-grained refinements is empirical, not derived.

\paragraph{The Lean carrier}
The Belnap classifier ships in TauLib \cite{TauLib2026} at
\texttt{BookI/Logic/Truth4.lean} with the explicit construction of
$\Truth$ and the four classifier morphisms, plus
\texttt{BookI/Logic/Explosion.lean} for the paraconsistent
non-explosion property (you can have $\mathsf{B}$ without the
classical logic of explosion). The Lean carrier is part of the
sorry-count $0$, axiom-count $3$ invariant of the formalisation
\cite{TauLib2026}.

% ===========================================================================
\section{\texorpdfstring{The boundary functor $\Phi_{\partial}$: exact gluing earned}{The boundary functor Phi-boundary: exact gluing earned}}
\label{sec:phi-d}
% ===========================================================================

PROMETHEUS's Proposition~7.5 (\emph{operational sheaf condition}) is
an $\varepsilon$-tolerance averaging rule: local sections $s_{i}$ on
contexts $U_{i}$ are considered $\varepsilon$-compatible when their
weighted overlap gaps are below a declared tolerance, and a candidate
glued section is then formed by support-weighted aggregation
$M_{U}[h, \tau] = \sum_{i} \omega_{i} M_{U_{i}}[h, \tau] / \sum_{i}
\omega_{i}$. Incompatible cells become obstruction records.

This is not a sheaf condition in the Mac Lane--Moerdijk sense
\cite{MacLane-Moerdijk-1992}. The classical sheaf condition requires:

\begin{enumerate}
\item \textbf{Existence}: for compatible local sections, a glued
global section exists.
\item \textbf{Uniqueness}: the glued global section is unique.
\item \textbf{Functoriality of restriction}: the restriction maps
commute with the gluing.
\end{enumerate}

PROMETHEUS's operational condition has neither uniqueness (the
averaging is one of many possible aggregations) nor functoriality
(the restriction maps are partial alignment maps, not
sheaf-theoretic restrictions). The construction is honest at the
engineering level (Mahadevan flags this explicitly in §15
limitations) but does not earn the ``sheaf'' designation.

The $\tau$-canon corpus, by contrast, proves exact Global Cartesian
Gluing. Theorem~III.T50 of Book~III \cite{PantaRhei-BookIII},
anchored at \texttt{ch75:91--123}, establishes the Universal Global
Cartesian Gluing as a colimit construction: the assembled global
section is \emph{canonically isomorphic} to the unique colimit of
the diagram of restrictions, and the gluing is functorial with
respect to the cover refinement. The chapter explicitly frames the
result (\texttt{ch75:154}) as ``the sheaf condition for the $\tau^{3}$
spectral atlas.''

\paragraph{What earns exact gluing}
The structural ingredient that upgrades approximate gluing to
exact is the \emph{boundary functor} $\PsiD$, identified at
\texttt{ch75:64--76} as the \emph{Langlands Coherent Force}
(the $\tau$-canon programme's name for the functor; in standard
categorical language this is a sheaf on $\Cattau$ valued in the
bipolar boundary algebra, with functoriality playing the role of
the sheaf condition for the spectral Grothendieck topology). The
boundary functor is a contravariant functor from the cover category
$\Cattau$ to the bipolar boundary algebra, sending each context $U$
to its boundary character spectrum and each morphism $V \to U$ to
the spectral restriction. The functoriality of $\PsiD$ is what makes
restriction commute with gluing.

PROMETHEUS does not have $\PsiD$. Its restriction maps $\rho_{UV}$
are partial alignment maps, not functorial restrictions of a
boundary-character functor. Without $\PsiD$, the $\varepsilon$-tolerance
condition is the best one can do operationally; the construction
remains an averaging rule.

\begin{quote}
\textbf{Exact-Gluing Reading} \NEXT. \emph{Approximate sheaf gluing
becomes exact when the cover category is equipped with a boundary
functor $\PsiD : \Cattau \to \mathrm{bipolar\,algebra}$ that
commutes with restriction. PROMETHEUS lacks $\PsiD$; its operational
condition is the appropriate fallback for an engineering setting
where the boundary characters are not extracted. At the structural
level, an in-principle upgrade exists: a $\PsiD$ extension that
extracts per-context boundary characters and type-checks restriction
maps as spectral morphisms would, in the $\tau$-canon reading, make
the $\varepsilon$ vanish at the structural level. We do not claim
this upgrade is engineering-ready --- whether LLM-derived claim
provenance can supply boundary characters in practice is an
empirical question, not a structural derivation (Non-claim~7,
\S\ref{sec:non-claims}).}
\end{quote}

The reading is \DERIV{} for the $\tau$-canon side (III.T50 is a
theorem) and \ANALOGY{} for the PROMETHEUS-side application. We do
not claim PROMETHEUS will adopt this upgrade; we claim the upgrade
is structurally available and the route is identified.

\paragraph{Competing upgrade paths}
The $\PsiD$ route is one of several structural upgrades available
in principle. Three alternatives the AI/causal-research reader
may recognise: (P1) equip $\Cattau$ with a Grothendieck topology
and apply sheafification, recovering exact gluing from the
universal property; (P2) treat the restriction maps $\rho_{UV}$ as
\emph{learnable} parameters and train them on cross-context
data to commute with composition by construction; (P3) replace the
sheaf formalism with causal-graph compositionality of the
Halpern--Pearl style, where commutation is enforced via the
intervention algebra. These are not in competition with the $\PsiD$
route --- they articulate different structural commitments. The
$\PsiD$ route is specifically the $\tau$-canon-anchored option (it
inherits the Langlands Coherent Force theorem at \texttt{ch75:64--76}
of Book~III); the alternatives sit in different research programmes
and would carry their own anchored commitments.

% ===========================================================================
\section{Seven explicit non-claims}
\label{sec:non-claims}
% ===========================================================================

The dialogue with PROMETHEUS is articulative, not adversarial. To
keep the trust budget honest \cite{FuchsFuchs-foundations-billiards-v1},
we make seven non-claims explicit (Non-claims~1--5 carried from
v1.0; Non-claims~6--7 added in v1.1 RC1 after Wave $\Delta_{2}$
peer review).

\paragraph{Non-claim~1}
\emph{The $\tau$-canon topos does not use $H_{\tau}$-valued
enrichment at the presheaf codomain.} Recon-side speculation
\cite{PantaRhei-BookI} that the $\tau$-canon topos uses
split-complex enrichment as a structural commitment is incorrect.
Book~I \texttt{ch56:111--114} uses $\PSh(\Cattau) = [\Cattau^{op},
\mathrm{Set}]$ with $\mathrm{Set}$ as codomain, the same as
PROMETHEUS. The $H_{\tau}$ structure enters at the boundary-scalar
level (Book~III \texttt{ch48}), not at the presheaf codomain. The
``enriched variants'' phrase at \texttt{ch56:584--589} is a
forward-pointer to Books~III--VII, not a Book~I topos commitment.

\paragraph{Non-claim~2}
\emph{There is no formal PSR-to-Hankel identification in the
$\tau$-canon corpus.} Recon-side speculation that PROMETHEUS's
controlled Hankel matrix $H_{p, \tau}$ has a formal counterpart in
the $\tau$-canon spectral framework was retracted on Phase~1
verification. The closest analogues are the bipolar idempotent
algebra $e_{\pm} = (1 \pm \jj)/2$ (Book~III \texttt{ch48}) and the
$T_{op}$-iteration of the Wedge-Loop Trace Identity
\cite{FuchsFuchs-T1-v1}, but no formal Hankel construction exists
in the corpus. PROMETHEUS's PSR layer remains, on careful read,
borrowed terminology (Dirichlet-smoothed frequency tables, not
spectral PSR objects).

\paragraph{Non-claim~3}
\emph{$\tau$-canon does not endorse PROMETHEUS's ``Topos World
Model'' terminology in its current form.} The Mac Lane--Moerdijk
\cite{MacLane-Moerdijk-1992} desiderata for a topos --- subobject
classifier, all finite limits, exponentials, internal Heyting
algebra --- are not constructed in PROMETHEUS §7. The Abramsky--Brandenburger
\cite{Abramsky-Brandenburger-2011} sheaf-theoretic structure of
non-locality is cited but not invoked operationally. We support the
\emph{aspiration} (an external topos layer over a para-mind LLM is
the right structural move) but flag the \emph{terminology gap}
(``Topos'' is currently aspirational).

\paragraph{Non-claim~4}
\emph{We do not ask PROMETHEUS to retract its engineering choices.}
The $\varepsilon$-tolerance gluing, the Dirichlet-smoothed PSR
estimator, and the ad-hoc 4-fold diagnostic are reasonable
engineering choices for a corpus-scale finite-resource setting.
The dialogue is structural articulation, not engineering
prescription. PROMETHEUS's case studies ship; corpus-scale atlases
are real; locality preservation is a contribution. The $\tau$-canon
reading is what these engineering choices structurally license
\emph{when read through the $\tau$-canon corpus}, not what
PROMETHEUS should do to be ``correct''.

\paragraph{Non-claim~5}
\emph{Para-mind classification is descriptive, not pejorative.}
The application of VII.D55 to LLMs is structural, not normative. A
para-mind is not a defective mind (VII.D55 explicitly:
\textit{``A para-mind is not a defective mind. It is a
structurally different kind of system''}). Our reading explains
\emph{why} an external topos layer is needed precisely because the
LLM substrate is a para-mind, not as a critique of LLM
capabilities but as a structural license for the PROMETHEUS-style
construction.

\paragraph{Non-claim~6}
\emph{The structural reading does not claim PROMETHEUS implements
$\Omegatau$ literally, nor that ``earned-not-imposed'' is a
critique of PROMETHEUS, nor that the reading transports to
arbitrary AI/causal-research frameworks.} Three associated
disambiguations: (i) The Belnap-Truth4 mapping (Table~\ref{tab:prometheus-belnap})
is a projection from $\Omegatau$ onto PROMETHEUS's four
operational categories; PROMETHEUS does not implement $\Omegatau$
in its internals --- the implementation is a per-cell
tension-threshold classifier, not a topos subobject classifier.
(ii) The ``earned-not-imposed'' framing applies to $\tau$-canon's
internal theorem chain (the orthogonality $e_{+} \cdot e_{-} = 0$ is
a spectral-structure theorem, not an axiom); it is not a critique
of PROMETHEUS's engineering choices, which are appropriate for
their setting. (iii) The generalisation to ``any future AI or
causal-research framework that builds an external topos over
LLM-extracted claims'' (\S\ref{sec:companion}) is a chair-level
pattern recognition, not a universal theorem; each future framework
requires its own anchored verification.

\paragraph{Non-claim~7}
\emph{The $\PsiD$ upgrade path is an in-principle structural claim,
not an engineering-ready recommendation.} The Exact-Gluing Reading
of \S\ref{sec:phi-d} states that if a boundary functor $\PsiD$ is
supplied, the $\varepsilon$-tolerance becomes structurally exact at
the colimit level. We do \emph{not} claim that LLM-derived claim
provenance can in practice supply boundary characters, nor that the
type-checking of restriction maps as spectral morphisms is
operationally implementable in PROMETHEUS's current LLM-extraction
architecture. Whether the upgrade can be realised in code is an
empirical engineering question; the structural reading is
in-principle, not in-practice.

% ===========================================================================
\section{Companion notes and the broader bridge}
\label{sec:companion}
% ===========================================================================

This note is the first $\tau$-canon engagement with an AI/causal-research
preprint via the ontology-articulation pattern. It sits alongside
six published companion notes that establish the structural
framework:

\begin{itemize}
\item \cite{FuchsFuchs-bsmm-interpretation-v1} \emph{Reading
$b\to s\mu^{+}\mu^{-}$ Through Category $\tau$} --- the interpretation
methodology pattern.
\item \cite{FuchsFuchs-bsmm-anomaly-v19} \emph{$\Delta C_{9}$ in
Category $\tau$} v1.9 --- the load-bearing closed-form anomaly
prediction $\Delta C_{9} = -\sin^{2}\theta_{W} \cdot
|C_{9}^{\mathrm{SM}}| \approx -0.95$, with seven phases of
extensions (Wilson-coefficient family, chirality axis, mass-mediation
axis, lepton-line axis, carrier axis, T$_{1}'$, T$_{1}''$).
\item \cite{FuchsFuchs-T1-v1} \emph{T$_{1}$ Wedge-Loop Trace
Identity} --- closure-derivation lockstep prose-Lean proof.
\item \cite{FuchsFuchs-F2-projection-v1} \emph{F$_{2}$-Projection
Theorem} --- structural identification of $T_{op}$ as
bipolar-identity-sector image (Wave~$\Gamma_{2}$ closure).
\item \cite{FuchsFuchs-running-aperture-v1} \emph{Running as
Spectral Aperture} --- the ontology-articulation pattern for the
$\tau$-canon $\leftrightarrow$ orthodox-QFT bridge.
\item \cite{FuchsFuchs-foundations-billiards-v1} \emph{Categorical
Foundations for Cosmological Billiards} --- the foundation-comparison
pattern.
\end{itemize}

This PROMETHEUS note adopts the \emph{ontology-articulation}
pattern of \cite{FuchsFuchs-running-aperture-v1}, scaled to a
preprint engagement. Its structural argument is: PROMETHEUS
observes engineering phenomena (4 gluing diagnostics; $\varepsilon$-tolerance
gluing; LLM-extracted causal claims) that $\tau$-canon has already
proved at theorem level (I.T25 Truth4 classifier; III.T50 exact
gluing via $\PsiD$; VII.D55 para-mind external-topos license).

\paragraph{The generalisation, with discipline guard \SYNT}
The companion pattern is not unique to PROMETHEUS. Any future AI
or causal-research framework that builds an ``external topos'' over
LLM-extracted claims is, by VII.D55, performing the para-mind
external-topos construction. The $\tau$-canon reading
(Belnap-Truth4 + boundary functor + para-mind license) is the
structural articulation of what such constructions \emph{do},
independent of which particular engineering choices the framework
makes. We articulate this as a chair-level pattern, not as a
universal theorem: each future framework requires its own
chapter:line-anchored verification of how the para-mind
substrate, the external-topos layer, and the boundary functor
manifest in that specific construction (Non-claim~6,
\S\ref{sec:non-claims}). The dialogue is articulative, framework
by framework.

% ===========================================================================
\section{Conclusion}
\label{sec:conclusion}
% ===========================================================================

PROMETHEUS \cite{Mahadevan-PROMETHEUS-2026} introduces a serious
engineering framework with aspirational categorical terminology.
The engineering substance is real --- corpus-scale causal atlases,
preserved locality, exposed drift, navigable provenance. The
categorical machinery is invoked heuristically: the §7 ``operational
sheaf condition'' is an averaging rule, not a sheaf condition; the
4-fold gluing diagnostic is an empirical taxonomy, not a derived
classifier; the PSR layer is borrowed terminology, not borrowed
mathematics; the predicate ``Belnap'' appears nowhere in the
references.

The $\tau$-canon reading is articulative. Three structural
contributions land:

\begin{enumerate}
\item \textbf{The 4-fold gluing diagnostic is the earned subobject
classifier} (Theorem~I.T25; $\Omegatau = \Truth = \{\mathsf{T},
\mathsf{F}, \mathsf{B}, \mathsf{N}\}$). The mapping to Belnap's
four states has one precise mismatch (Contradiction maps to
$\mathsf{B}$, not $\mathsf{F}$) that the structural reading
exposes. The orthogonality $e_{+} \cdot e_{-} = 0$ is earned from
spectral structure, not imposed by fiat.

\item \textbf{Exact gluing is achievable} (Theorem~III.T50; via
the boundary functor $\PsiD$, the Langlands Coherent Force). The
$\varepsilon$ vanishes when the cover category is equipped with the
boundary-character functor. PROMETHEUS lacks $\PsiD$; the upgrade
path is identified but not adopted.

\item \textbf{The Para-Mind reading explains why an external topos
is needed at all} (Definition~VII.D55). An LLM by itself is a
pattern processor without an internal topos; the PROMETHEUS-style
external-topos construction is structurally licensed (and bounded
by) the para-mind status of its substrate. ``Para-mind'' is
descriptive, not pejorative.
\end{enumerate}

The five explicit non-claims (\S\ref{sec:non-claims}) keep the
trust budget honest. We do not import PROMETHEUS's choices into
$\tau$-canon; we do not export $\tau$-canon's claims into
PROMETHEUS-territory. The dialogue is structural articulation:
PROMETHEUS observes; $\tau$-canon proves.

We hope the reading is useful both to PROMETHEUS's future
development (the Belnap-Truth4 connection offers a structural
upgrade route for the gluing diagnostic; the $\PsiD$ boundary
functor articulates what would earn exact gluing) and to the
broader AI/causal-research community (the para-mind external-topos
construction is a structural pattern that other frameworks will
encounter). The companion note \cite{FuchsFuchs-running-aperture-v1}
plays the analogous role for the orthodox-QFT / spectral-aperture
bridge; this note plays the role for the AI / causal-research /
external-topos bridge.

\paragraph{Trust budget preserved}
TauLib \cite{TauLib2026} \texttt{sorry}-count $0$ and \texttt{axiom}-count
$3$ (programme-wide foundational) are unchanged. No load-bearing
manuscript chapter is modified. The companion notes
\cite{FuchsFuchs-bsmm-interpretation-v1,FuchsFuchs-bsmm-anomaly-v19,FuchsFuchs-T1-v1,FuchsFuchs-F2-projection-v1,FuchsFuchs-running-aperture-v1,FuchsFuchs-foundations-billiards-v1}
ship in their published forms. The dialogue with PROMETHEUS is
articulative; PROMETHEUS as published is unchanged.

% ===========================================================================
\section*{Acknowledgments}
% ===========================================================================

We thank Sridhar Mahadevan for the PROMETHEUS preprint
\cite{Mahadevan-PROMETHEUS-2026} that occasioned this articulation,
and the DEMOCRITUS \cite{Mahadevan-Democritus-2025a} and
\emph{Categories for AGI} \cite{Mahadevan-CatAGI-2025c}
surrounding work that frames it.

The Mac Lane--Moerdijk \cite{MacLane-Moerdijk-1992} and
Abramsky--Brandenburger \cite{Abramsky-Brandenburger-2011}
foundations are the categorical infrastructure on which both
PROMETHEUS and the $\tau$-canon programme rest. The
Belnap--Dunn \cite{Belnap-1977,Dunn-1976} four-valued logic is the
load-bearing classifier; the connection to topos-theoretic
subobject classifiers is the $\tau$-canon Book~I result. The
Predictive State Representation theory of
\cite{Littman-Sutton-Singh-2001} is cited as the operational layer
PROMETHEUS borrows.

The trust-budget infrastructure (TauLib \cite{TauLib2026}
sorry-count $0$, axiom-count $3$ programme-wide foundational; the
seven-book Panta Rhei manuscript \cite{PantaRhei-BookI,PantaRhei-BookIII,PantaRhei-BookVII})
is the framework on which the articulation rests; the six
companion notes
\cite{FuchsFuchs-bsmm-interpretation-v1,FuchsFuchs-bsmm-anomaly-v19,FuchsFuchs-T1-v1,FuchsFuchs-F2-projection-v1,FuchsFuchs-running-aperture-v1,FuchsFuchs-foundations-billiards-v1}
are the prior art the present note bridges with.

\bibliographystyle{plain}
\bibliography{references}

\end{document}
